% Chapter 1, typeset from lectures/L01: the README is the opener and the closing
% review, and each appendix is a section, in the same order and with the same subsections.

\chapter{Zephyr, west and the Build}
\label{c1:ch:build}\label{c1:ch}

{\large\itshape What Zephyr is, and how an application becomes a program.\par}

\begin{chapterbox}{In this chapter}
\begin{itemize}
  \item What Zephyr is: a kernel, drivers, subsystems, and a build system that ties them together.
  \item \code{native_sim}: Zephyr as a Linux program, and why this book runs on it.
  \item The west workspace: the manifest, \code{west update}, and pinning a Zephyr release.
  \item Anatomy of an application: \code{CMakeLists.txt}, \code{prj.conf} and \code{src/}.
  \item Building and running: \code{west build}, the build directory, and \code{west build -t run}.
  \item Kconfig: options, dependencies, where values come from, and how to read what the build chose.
  \item Application-specific Kconfig options, used from C and from CMake.
  \item \code{printk()} and the logging subsystem: modules, levels and compile-time filtering.
\end{itemize}
\end{chapterbox}

This chapter is about everything that happens before a line of your code is compiled. Zephyr decides
what a build contains, and what the hardware in it is, long before the compiler runs, and nearly
every question a newcomer has is really a question about that: where Zephyr lives, what west does
with it, which options a build ended up with, and why an option that was set was not the one used.

Work the chapter in order. \Secref{c1:app:a} sets up the workspace: do it before the first lecture,
because fetching Zephyr takes a few minutes. \Secref{c1:app:b} is the substance of the chapter, and
the exercises in \secref{c1:sec:exercises} are where you write the code yourself. The examples the
lecture is taught from are \file{lectures/L01/examples/hello} and
\file{lectures/L01/examples/kconfig} in the companion repository.

% Section 1.1, from lectures/L01/appendix/a_setup.md.

\section{Setting Up the Workspace}
\label{c1:sec:settinguptheworksp}\label{c1:app:a}
Every program in this book is built and run on Linux. This section installs what the workspace
needs, sets it up, and checks that it works. Do it before the first lecture.

\subsection{Linux}
Any recent x86-64 Linux works. On Windows, use \textbf{WSL 2} with \textbf{Ubuntu 24.04}, as in the
other courses in this track. In \textbf{Windows PowerShell as Administrator}:

\begin{shell}
wsl --install -d Ubuntu-24.04
\end{shell}

\noindent Restart if asked, start \emph{Ubuntu} from the Start menu, and choose a user name and
password. \code{native_sim} needs WSL \textbf{2}: the old WSL 1 cannot run the programs it builds.
Check with \code{wsl --status} in PowerShell, which must say \code{Default Version: 2}.

\textbf{Keep the repository in the Linux file system}, for example in \code{~/qacademy}, and not
under \code{/mnt/c}. A Zephyr build touches thousands of files, and through \code{/mnt/c} it is many
times slower.

\subsection{Packages}
In the Ubuntu terminal:

\begin{shell}
sudo apt -y update
sudo apt -y install git make cmake gcc g++ python3 python3-venv python3-pip clang-format gdb
\end{shell}

\noindent What each one is for:

\begin{booktable}{@{}LL@{}}
  \thead{Package} & \thead{Used for} \\
  \midrule
  \code{git} & Fetching the companion repository, and Zephyr through west. \\
  \code{make} & The repository’s \code{make setup}, \code{make test} and friends. \\
  \code{cmake} & Zephyr’s build system. Version 3.20 or newer; Ubuntu 22.04 and 24.04 both have it. \\
  \code{gcc}, \code{g++} & The compiler. \code{native_sim} programs are Linux programs, built by the host’s own GCC. \\
  \code{python3}, \code{python3-venv}, \code{python3-pip} & west, Twister, and the scripts that process Kconfig and devicetree. \\
  \code{clang-format} & \code{make format}, which formats C files the way Zephyr does. \\
  \code{gdb} & Debugging in Chapter~\ref{c6:ch}. \\
\end{booktable}

\textbf{No Zephyr SDK is needed.} The SDK holds cross-compilers for real microcontrollers;
\code{native_sim} uses the host compiler, which \code{source env.sh} (below) tells Zephyr to do.

Two packages are optional:

\begin{itemize}
  \item \code{device-tree-compiler} (\code{dtc}). Zephyr processes devicetree with its own Python scripts
    and uses \code{dtc} only for extra warnings. Without it, every build prints \code{Could NOT find Dtc
    (missing: DTC)}, which is harmless.
  \item \code{gcc-multilib g++-multilib}. This book builds for \code{native_sim/native/64}, a 64-bit
    program. The plain \code{native_sim} board is 32-bit, and building for it without these packages
    fails with \code{fatal error: bits/libc-header-start.h: No such file or directory}.
\end{itemize}

\subsection{The workspace}
Clone the repository and set it up:

\begin{shell}
git clone https://github.com/qrtech-academy/zephyr.git
cd zephyr
make setup
\end{shell}

\noindent \code{make setup} does four things, all inside the repository:

\begin{enumerate}
  \item Creates a Python environment in \code{.venv} and installs west into it.
  \item Makes the repository a west workspace, with \file{west.yml} as its manifest.
  \item Runs \code{west update}, which fetches the Zephyr release that \code{west.yml} pins into
    \code{deps/zephyr}.
  \item Installs Zephyr’s Python requirements, Twister’s among them, into \code{.venv}.
\end{enumerate}

It takes a few minutes and about 1 GB of disk. Nothing is installed outside the repository, and
running it again is safe.

Then, \textbf{in every new terminal}:

\begin{shell}
source env.sh
\end{shell}

\noindent That activates the Python environment, so that \code{west} is found, and sets two
variables: \code{ZEPHYR_BASE}, which says where Zephyr is, and \code{ZEPHYR_TOOLCHAIN_VARIANT=host},
which says to build with the host’s GCC. Your prompt starts with \code{(.venv)} while it is active.

\subsection{Checking that it works}
From the repository root:

\begin{shell}
west build -p -b native_sim/native/64 lectures/L01/examples/hello
west build -t run
\end{shell}

\noindent The first command configures and builds the example into \code{build/}, which takes about
half a minute the first time. The second runs it:

\begin{console}
*** Booting Zephyr OS build v4.4.2 ***
Hello from native_sim/native/64
[00:00:00.000,000] <inf> hello: Tick 1 at 0 ms
[00:00:00.510,000] <inf> hello: Tick 2 at 510 ms
[00:00:01.020,000] <inf> hello: Tick 3 at 1020 ms
[00:00:01.530,000] <wrn> hello: A warning, in yellow
\end{console}

\noindent The program keeps running after that, because a Zephyr system does not exit when
\code{main()} returns. Stop it with \code{Ctrl+C}. Finally, check that the test suites run:

\begin{shell}
make test
\end{shell}

\noindent With no exercises written yet, every exercise set is reported as \code{SKIP}, and the
command succeeds.

\subsection{Visual Studio Code}
With the \textbf{WSL} extension, run \code{code .} from the repository root in the Ubuntu terminal.
For code navigation that understands Zephyr’s include paths and \code{CONFIG_} macros, install the
\textbf{C/C++} extension and point it at the build’s compilation database, which every build writes
to \code{build/compile_commands.json}. In \code{.vscode/settings.json}:

\begin{yamlcode}
{
    "C_Cpp.default.compileCommands": "${workspaceFolder}/build/compile_commands.json"
}
\end{yamlcode}

\subsection{When something goes wrong}
\begin{booktable}{@{}LL@{}}
  \thead{Message} & \thead{Cause} \\
  \midrule
  \code{west: command not found} & \code{source env.sh} has not been run in this terminal. \\
  \code{ZEPHYR_TOOLCHAIN_VARIANT not set, trying to locate Zephyr SDK} followed by \code{Could not find a package configuration file provided by "Zephyr-sdk"} & Also \code{source env.sh}: without it, Zephyr looks for the SDK. \\
  \code{error: the workspace is not set up} & \code{make setup} has not been run, or did not finish. Run it again. \\
  \code{fatal error: bits/libc-header-start.h: No such file or directory} & Building for the 32-bit \code{native_sim} without \code{gcc-multilib}. Use \code{native_sim/native/64}. \\
  A build that fails for no reason you can see after changing boards or options & A stale build directory. Build again with \code{-p} (pristine). \\
\end{booktable}

% Section 1.2, from lectures/L01/appendix/b_zephyr_and_the_build.md.

\section{Zephyr, west and the Build}
\label{c1:sec:zephyrwestandthebu}\label{c1:app:b}
This section covers what Zephyr is, the board this book runs on, the workspace it lives in, and the
build that turns an application into a program. By the end you should be able to explain every file
in the \file{lectures/L01/examples/hello} and \file{lectures/L01/examples/kconfig} examples, and
every line of output their builds print that matters.

\subsection{What Zephyr is}
Zephyr is an open-source real-time operating system for microcontrollers, hosted by the Linux
Foundation and licensed under Apache 2.0. It runs on ARM Cortex-M and Cortex-A, RISC-V, x86, Xtensa,
ARC and more, and ships support for well over a thousand boards.

The word \emph{operating system} is meant more literally than in ``FreeRTOS''. FreeRTOS gives you a
scheduler and a handful of kernel objects, and you bring everything else: the drivers, the logging,
the build. Zephyr ships all of it:

\begin{booktable}{@{}>{\raggedright\arraybackslash}p{10.5em}L@{}}
  \thead{Part} & \thead{What it gives you} \\
  \midrule
  \textbf{The kernel} & Threads, scheduling, interrupts, timers, and kernel objects for synchronization and communication. Chapters~\ref{c3:ch} and~\ref{c4:ch}. \\
  \textbf{The device driver model} & One API per kind of device (GPIO, I2C, sensor, flash, …), implemented by drivers for many chips. Chapters~\ref{c2:ch} and~\ref{c5:ch}. \\
  \textbf{Subsystems} & Logging, a shell, settings storage, file systems, networking, Bluetooth, USB, and more. This chapter and Chapter~\ref{c6:ch}. \\
  \textbf{The build system} & CMake, Kconfig and devicetree, which select and configure all of the above at compile time. This chapter and Chapter~\ref{c2:ch}. \\
  \textbf{Test tooling} & The ztest framework and the Twister test runner. Chapter~\ref{c5:ch}. \\
\end{booktable}

\textbf{Everything is decided at compile time.} A Zephyr image holds only the drivers, subsystems
and kernel features the application asked for, sized for the board it was built for. There is no
loading of drivers at run time and nothing to discover: the build has already worked out what
hardware exists and which code serves it. That is what makes Zephyr fit into a few kilobytes, and it
is also why most of what is hard about Zephyr is in the build. If you have met Linux’s \code{make
menuconfig} or its \code{.dts} files, those are the same two ideas, Kconfig and devicetree, used the
same way.

Zephyr releases roughly three times a year, with a long-term support (LTS) release, maintained with
fixes for years, every few releases. Projects pin a release and move deliberately; so does this
book.

\subsection{\code{native_sim}}
A \emph{board}, in Zephyr, is what an application is built for: a processor, the peripherals around
it, and how they are wired. \code{native_sim} is a board that does not exist. It is Zephyr’s
\textbf{POSIX architecture}: the kernel, drivers and application compiled together into an ordinary
Linux program, which runs on the host CPU.

\begin{shell}
west build -b native_sim/native/64 lectures/L01/examples/hello
./build/zephyr/zephyr.exe
\end{shell}

Nothing is emulated at the instruction level, as QEMU would. Each Zephyr thread is backed by a Linux
thread, and the kernel lets exactly one of them run at a time, so Zephyr’s scheduler is in charge
exactly as on a microcontroller. The hardware is what is simulated: a system timer, an interrupt
controller, and a set of peripherals with emulated backends, among them a GPIO controller, an I2C
bus, a flash chip and an EEPROM. Chapters~\ref{c2:ch} and~\ref{c5:ch} use those.

\textbf{Time on \code{native_sim} is simulated.} The system timer counts simulated microseconds, and
the simulator advances that clock only when every Zephyr thread is waiting for something. By default
a program slows its clock down to real time, so that a one-second sleep takes a second, which is
what you want when you watch it. Two command-line options change that:

\begin{booktable}{@{}lL@{}}
  \thead{Option} & \thead{Effect} \\
  \midrule
  \code{-no-rt} & Advance simulated time as fast as the host can. Ten simulated minutes take a moment. \\
  \code{-stop_at=<seconds>} & Stop the program after that much simulated time. \\
\end{booktable}

\code{./build/zephyr/zephyr.exe --help} lists all of them. Because time is simulated, a program’s
behaviour does not depend on how busy the host is: it prints the same timestamps on every run, on
every machine. The exercises rely on that.

\textbf{The same property is \code{native_sim}’s biggest limitation.} Code runs in zero simulated
time, so no interrupt, not even the timer’s, can arrive while a thread is running; only when it
waits. A loop such as \code{while (!flag) \{\}}, waiting for an interrupt to set \code{flag}, never
ends, and because time cannot advance, nothing else runs either. On a microcontroller that loop is
merely wasteful. Chapter~\ref{c3:ch} shows it on purpose.

A few kernel features are also missing on \code{native_sim}: memory protection between threads
(userspace), and the stack overflow checks, since each thread really runs on its Linux thread’s
stack. The Zephyr documentation lists them under \emph{Native simulator: Important limitations}.

\subsection{The west workspace}
Zephyr is not one repository. The kernel and core live in \code{zephyrproject-rtos/zephyr}, and the
vendor HALs, cryptography libraries, file systems and other \emph{modules} each live in repositories
of their own, about eighty in all. \textbf{west} is Zephyr’s tool for managing that: a
multi-repository manager in the spirit of Git submodules or Google’s \code{repo}, plus a set of
commands for building, flashing and testing.

\subsubsection{Workspaces and manifests}
A \textbf{west workspace} is a directory with a \code{.west/config} file in it. The config names a
\textbf{manifest}, a \code{west.yml} file listing the Git repositories (\emph{projects}) that belong
to the workspace, where each one goes, and which revision of each one to check out.

The companion repository is its own workspace. \code{make setup} writes a \code{.west/config} that
points at the repository’s \file{west.yml}:

\begin{yamlcode}
manifest:
  remotes:
    - name: zephyrproject-rtos
      url-base: https://github.com/zephyrproject-rtos

  projects:
    - name: zephyr
      remote: zephyrproject-rtos
      revision: v4.4.2
      path: deps/zephyr
      clone-depth: 1
      west-commands: scripts/west-commands.yml

  self:
    path: .
\end{yamlcode}

\noindent Line by line:

\begin{itemize}
  \item \textbf{\code{remotes}} give a name to a base URL. A project’s repository URL is its remote’s
    \code{url-base} followed by the project’s name: \code{https://github.com/zephyrproject-rtos/zephyr}.
  \item \textbf{\code{revision: v4.4.2}} pins Zephyr to a release tag. It could as well be a branch or a
    commit hash, and in a product it should be a tag or a hash, so that the same manifest checks out the
    same code a year later.
  \item \textbf{\code{path: deps/zephyr}} is where the project is checked out, relative to the workspace.
  \item \textbf{\code{clone-depth: 1}} fetches only the pinned commit, not Zephyr’s decade of history, which
    makes the first \code{west update} much faster, at the cost of \code{git log} showing one commit.
  \item \textbf{\code{west-commands}} tells west that Zephyr provides extra commands, \code{west build} and
    \code{west twister} among them. Without that line, west knows only its built-in commands.
  \item \textbf{\code{self: path: .}} says that the manifest repository is the workspace directory itself.
\end{itemize}

What the manifest does \textbf{not} do is import Zephyr’s own \code{west.yml}, which lists all of
Zephyr’s modules. \code{native_sim} needs none of them, and together they are several gigabytes. A
product built for a real microcontroller would import the few it needs, such as the vendor’s HAL, by
name.

\subsubsection{The commands}
West’s own commands manage the workspace:

\begin{shell}
west update                  # Check out every project at the revision the manifest pins.
west list                    # Show every project, its path and its revision.
west manifest --resolve      # Show the whole manifest, with imports expanded.
west config build.board      # Read a configuration value from .west/config.
\end{shell}

\noindent \code{west list} in this workspace:

\begin{console}
manifest     .                            HEAD                                     N/A
zephyr       deps/zephyr                  v4.4.2                                   https://github.com/zephyrproject-rtos/zephyr
\end{console}

\noindent \textbf{Moving to another Zephyr release} is two steps: change \code{revision} in
\code{west.yml}, and run \code{west update}. West never updates a project on its own; \code{git
pull} in the companion repository changes the manifest, and \code{west update} is what makes the
workspace match it.

The extension commands that Zephyr adds are the ones used every day: \code{west build}, and, in
Chapter~\ref{c5:ch}, \code{west twister}. On real hardware there would also be \code{west flash} and
\code{west debug}.

\subsection{Anatomy of an application}
A Zephyr application is a directory with at least three files:

\begin{console}
hello/
    CMakeLists.txt      How to build it.
    prj.conf            How to configure Zephyr for it.
    src/main.c          The code.
\end{console}

\subsubsection{\code{CMakeLists.txt}}
\begin{cmakecode}
cmake_minimum_required(VERSION 3.20.0)
find_package(Zephyr REQUIRED HINTS $ENV{ZEPHYR_BASE})
project(hello)

target_sources(app PRIVATE src/main.c)
\end{cmakecode}

\noindent \textbf{\code{find_package(Zephyr)} is where Zephyr’s build system takes over.} It finds
Zephyr through \code{ZEPHYR_BASE}, which \code{env.sh} sets, and then does most of the work
described in \emph{The build} below: it reads the board, processes the devicetree, runs Kconfig, and
sets up the whole of Zephyr as a set of CMake libraries. It also creates a library target named
\textbf{\code{app}}, which is where the application’s own sources go. \code{project()} must come
after it, because Zephyr selects the compiler before CMake’s \code{project()} looks for one.

\subsubsection{\code{prj.conf}}
A list of Kconfig settings, one per line, applied on top of the board’s defaults:

\begin{kconfigcode}
CONFIG_LOG=y
\end{kconfigcode}

\noindent The file must exist, even if it is empty. An application without one fails to configure
with \code{No prj.conf file(s) was found}.

\subsubsection{\code{src/main.c}}
\code{main()} is called by the kernel after it has initialized itself and every driver, in a thread
called \code{main}. Returning from it ends that thread, \textbf{not the program}: the kernel keeps
running, along with every other thread, and if nothing else has work to do, the idle thread runs
forever. On \code{native_sim} that is why a program whose \code{main()} has returned keeps running
until \code{Ctrl+C} or \code{-stop_at}.

\subsubsection{Optional files}
Applications grow more files as they need them. Those you will meet in this book:

\begin{booktable}{@{}LLl@{}}
  \thead{File} & \thead{What it is} & \thead{Chapter} \\
  \midrule
  \code{Kconfig} & Configuration options of the application’s own. & This one \\
  \code{*.conf} & Extra Kconfig settings, merged when asked for with \code{-DEXTRA_CONF_FILE}. & This one \\
  \code{app.overlay} & Changes to the board’s devicetree for this application. & Chapter~\ref{c2:ch} \\
  \code{dts/bindings/} & Devicetree bindings of the application’s own. & Chapter~\ref{c2:ch} \\
  \code{boards/} & Kconfig settings and overlays that apply to one board only. & Chapter~\ref{c2:ch} \\
  \code{sample.yaml}, \code{testcase.yaml} & How Twister builds, runs and checks the application. & Chapter~\ref{c5:ch} \\
\end{booktable}

\subsection{The build}

\subsubsection{\code{west build}}
\begin{shell}
west build -b native_sim/native/64 lectures/L01/examples/hello
\end{shell}

\code{west build} is a wrapper around CMake and Ninja. It configures the application for the board
given with \code{-b}, into a build directory, \code{build/} unless \code{-d} says otherwise, and
then builds it. The options you will use most:

\begin{booktable}{@{}lL@{}}
  \thead{Option} & \thead{Effect} \\
  \midrule
  \code{-b <board>} & The board to build for. The repository sets \code{native_sim/native/64} as the workspace default, so it may be left out. \\
  \code{-d <dir>} & The build directory. Defaults to \code{build}. \\
  \code{-p} & \emph{Pristine}: delete the build directory first. By default (\code{-p auto}) west does so only when the application directory has changed. \\
  \code{-t <target>} & Build a particular target, such as \code{run} or \code{menuconfig}. \\
  \code{-- <args>} & Pass the rest on to CMake, e.g. \code{-- -DEXTRA_CONF_FILE=shout.conf}. \\
\end{booktable}

A build directory remembers its board and application, so a second build of the same application
needs no arguments, and \code{west build -t run} from the repository root runs what \code{build/}
holds.

\subsubsection{What happens when}
A Zephyr build has two phases, and the first is the one that is different.

\textbf{Configuration}, which CMake runs, happens once per build directory, and again when an input
changes. In order:

\begin{enumerate}
  \item \textbf{The board} is found, and with it the processor architecture, the SoC, and the board’s
    default Kconfig settings and devicetree.
  \item \textbf{Devicetree.} The board’s devicetree and any overlays are combined into one tree, checked
    against their bindings, and turned into a C header of macros.
  \item \textbf{Kconfig.} Every Kconfig option in Zephyr and the application is resolved, from its defaults,
    the board’s settings, \code{prj.conf} and the rest, into one final configuration. Kconfig runs after
    devicetree because options may depend on it: once GPIO is enabled at all, the driver for a GPIO
    controller is enabled by default exactly when the devicetree has an enabled node for it
    (\code{depends on DT_HAS_ZEPHYR_GPIO_EMUL_ENABLED}, for \code{native_sim}’s).
  \item \textbf{CMake} generates the build: which of Zephyr’s source files to compile, decided by the
    Kconfig options, with which flags.
\end{enumerate}

\textbf{Compilation}, which Ninja runs, then builds the libraries Zephyr is made of and links them
with the application into one image. On \code{native_sim} that image is a Linux executable.

\bookfigure{build_pipeline}{The build, in the order it happens. Configuration runs once per build
directory, and again when an input changes; compilation follows it. Kconfig comes after devicetree,
so that an option can depend on a node.}{c1:fig:build}

The order explains most of Zephyr’s error messages. A Kconfig error aborts configuration before a
line of C is compiled. A devicetree error aborts it even earlier. And a C file cannot change which
drivers are built, because by the time it is compiled, that has long been decided.

\subsubsection{The build directory}
Among the thousands of files in \code{build/}, a few are worth knowing by name:

\begin{booktable}{@{}LL@{}}
  \thead{File} & \thead{What it holds} \\
  \midrule
  \code{build/zephyr/zephyr.exe} & The program, on \code{native_sim}. On a microcontroller the image would be \code{zephyr.elf}, \code{zephyr.bin} and \code{zephyr.hex}; \code{native_sim} writes a \code{zephyr.elf} too. \\
  \code{build/zephyr/.config} & The final Kconfig configuration: every option and its value. \\
  \code{build/zephyr/include/generated/zephyr/autoconf.h} & The same configuration as C macros, included into every source file. \\
  \code{build/zephyr/zephyr.dts} & The final devicetree, overlays applied. Chapter~\ref{c2:ch}. \\
  \code{build/zephyr/include/generated/zephyr/devicetree_generated.h} & The devicetree as C macros. Chapter~\ref{c2:ch}. \\
  \code{build/zephyr/zephyr.map} & The linker map: where every symbol went. \\
  \code{build/compile_commands.json} & Every compiler command line, for editors. \\
\end{booktable}

\subsubsection{Running}
On \code{native_sim}, the program is run either by the build or directly:

\begin{shell}
west build -t run                             # Run build/zephyr/zephyr.exe.
./build/zephyr/zephyr.exe -stop_at=3 -no-rt   # Run it with options.
\end{shell}

\subsection{Kconfig}
Zephyr has thousands of configuration options. Kconfig, the configuration language from the Linux
kernel, is how they are declared, how their dependencies are enforced, and how an application sets
them.

\subsubsection{Declaring options}
An option is a \emph{symbol} with a type, a prompt, a default, and optionally dependencies:

\begin{kconfigcode}
config LOG
	bool "Logging"
	select PRINTK if USERSPACE
	help
	  Global switch for the logger, when turned off log calls will not be
	  compiled in.

if LOG

config LOG_DEFAULT_LEVEL
	int "Default log level"
	default 3
	range 0 4

endif # LOG
\end{kconfigcode}

\noindent (Both are Zephyr’s own, from \code{subsys/logging/Kconfig} and \code{Kconfig.filtering},
shortened.)

\begin{booktable}{@{}lL@{}}
  \thead{Keyword} & \thead{Meaning} \\
  \midrule
  \code{bool}, \code{int}, \code{hex}, \code{string} & The type, followed by the prompt shown in menuconfig. A symbol with no prompt cannot be set by the user; only defaults and \code{select} set it. \\
  \code{default} & The value when nothing sets it. Several \code{default ... if ...} lines may follow each other; the first whose condition holds applies. \\
  \code{depends on} & The symbol does not exist, and is \code{n}, unless the condition holds. \code{if LOG} … \code{endif} adds \code{depends on LOG} to every symbol between them. \\
  \code{select} & Forces another bool symbol to \code{y} whenever this one is \code{y}, dependencies or not. \\
  \code{range} & The allowed values of an \code{int} or \code{hex}. \\
  \code{menu}, \code{choice} & Group symbols in menuconfig; a \code{choice} allows exactly one of its bool symbols. \\
\end{booktable}

\subsubsection{Where a value comes from}
An option’s final value is the last one assigned, from these sources in this order:

\begin{enumerate}
  \item The \textbf{default} in the option’s declaration.
  \item The \textbf{board’s defaults}, from the board’s \code{_defconfig} file. For this book’s board, that
    is \code{deps/zephyr/boards/native/native_sim/native_sim_64_defconfig}.
  \item \textbf{\code{prj.conf}}, and any board-specific \code{.conf} file in the application’s
    \code{boards/} directory.
  \item \textbf{Extra configuration files}, given with \code{-DEXTRA_CONF_FILE=<file>}.
  \item \textbf{The command line}: \code{-DCONFIG_LOG_DEFAULT_LEVEL=2} after \code{--} on the \code{west
    build} line.
\end{enumerate}

Then Kconfig checks the result. Every assignment is only a \emph{request}: an option whose
dependencies do not hold stays \code{n}, and an \code{int} outside its range falls back on its
default. \textbf{Zephyr treats that as an error}, because a request that silently did not take is
how configurations rot. With \code{CONFIG_SHELL_LOG_BACKEND=y} but no shell:

\begin{console}
warning: SHELL_LOG_BACKEND (defined at subsys/shell/Kconfig:278) was assigned the value 'y' but got
the value 'n'. Check these unsatisfied dependencies: SHELL (=n). ...

error: Aborting due to Kconfig warnings
\end{console}

\noindent The message names the dependency that failed. The fix is to enable it too, or to find out
why it cannot be.

\subsubsection{Reading the result}
After a build, \code{build/zephyr/.config} holds every symbol’s value, and grep is the fastest way
to answer \emph{what did the build actually use?}:

\begin{shell}
grep -E "CONFIG_LOG_MODE|CONFIG_SYS_CLOCK_TICKS_PER_SEC" build/zephyr/.config
\end{shell}

\begin{kconfigcode}
CONFIG_SYS_CLOCK_TICKS_PER_SEC=100
CONFIG_LOG_MODE_IMMEDIATE=y
\end{kconfigcode}

\noindent For browsing, \code{west build -t menuconfig} opens the same configuration in the
terminal, and \code{west build -t guiconfig} in a window, if Ubuntu’s \code{python3-tk} package is
installed. \code{/} searches, and \code{?} on a symbol shows its definition, dependencies and what
selects it, which is the fastest way to find out \emph{why} an option has the value it has.
\textbf{Changes made in menuconfig live only in the build directory}: they are lost with the next
pristine build. Anything that should last belongs in \code{prj.conf}.

\subsubsection{Using options from C}
Every option that is set becomes a macro in \code{autoconf.h}, which the build includes into every
source file: \code{CONFIG_LOG_DEFAULT_LEVEL} is \code{3}, \code{CONFIG_BOARD_TARGET} is
\code{"native_sim/native/64"}, and a bool that is \code{y} is \code{1}. \textbf{A bool that is
\code{n} is not defined at all.} There are two ways to test one:

\begin{ccode}
#ifdef CONFIG_APP_SHOUT
	shout(greeting, loud, sizeof(loud));
#endif

if (IS_ENABLED(CONFIG_APP_SHOUT)) {
	shout(greeting, loud, sizeof(loud));
}
\end{ccode}

\noindent Prefer \textbf{\code{IS_ENABLED()}}. It expands to \code{1} or \code{0}, so the code
inside is compiled and type-checked in every configuration and removed by the compiler in the ones
where it cannot run. Code behind \code{\#ifdef} is not even parsed when the option is off, so it
breaks unnoticed. Keep \code{\#ifdef} for what cannot compile at all when the option is off, such as
a reference to a struct member that exists only with the option.

Both forms share one trap. \code{IS_ENABLED(CONFIG_APP_SHOUTT)}, misspelled, compiles without a
warning and is \code{0}, because an undefined macro is simply an option that is off. Kconfig catches
a misspelled option in \code{prj.conf}; nothing catches one in C.

\subsubsection{Using options from CMake}
Kconfig has run by the time \code{find_package(Zephyr)} returns, so every option is a CMake variable
too, and CMake can decide what to compile:

\begin{cmakecode}
target_sources_ifdef(CONFIG_APP_SHOUT app PRIVATE src/shout.c)
\end{cmakecode}

\noindent That is how all of Zephyr is built: nearly every driver and subsystem is a list of
\code{*_sources_ifdef(CONFIG_...)} lines.

\subsubsection{Options of your own}
An application with a file named \code{Kconfig} gets options of its own. That file \textbf{replaces}
Zephyr’s top-level Kconfig file for the build, so it must bring Zephyr’s back in at the end:

\begin{kconfigcode}
menu "Application"

config APP_REPEAT
	int "How many times to print the greeting"
	range 1 10
	default 3

endmenu

source "Kconfig.zephyr"
\end{kconfigcode}

\noindent Prefix the options with \code{APP_}, or with the application’s name, so that they cannot
collide with Zephyr’s. They then behave exactly like Zephyr’s own: they appear in menuconfig, are
set from \code{prj.conf}, and reach C and CMake as \code{CONFIG_APP_REPEAT}.

\subsection{\code{printk} and logging}

\subsubsection{\code{printk}}
\code{printk()} prints to the console, with \code{printf()}-style formatting from Zephyr’s own small
formatter. It writes synchronously, from any context, and is on by default. It is the right tool for
a quick look, and the wrong one for anything that ships.

\subsubsection{The logging subsystem}
Logging is enabled with \code{CONFIG_LOG=y}. Each source file belongs to a \textbf{log module},
registered once, with a name and a level:

\begin{ccode}
#include <zephyr/logging/log.h>

LOG_MODULE_REGISTER(monitor, LOG_LEVEL_INF);

LOG_ERR("Temperature %d C is critical", celsius);
LOG_WRN("Temperature %d C is high", celsius);
LOG_INF("Temperature %d C", celsius);
LOG_DBG("Raw sample 0x%04x", raw);
\end{ccode}

\noindent Every message is printed with a timestamp, its level and its module, and in color on a
terminal:

\begin{console}
[00:00:01.530,000] <wrn> monitor: Temperature 33 C is high
\end{console}

\noindent Other files of the same module use \code{LOG_MODULE_DECLARE(monitor, LOG_LEVEL_INF)}
instead of registering it again. \code{LOG_HEXDUMP_INF(data, length, "label")} prints a buffer.

\subsubsection{Levels}
The levels are numbered: \code{LOG_LEVEL_NONE} is 0, \code{ERR} 1, \code{WRN} 2, \code{INF} 3 and
\code{DBG} 4. \textbf{A module’s level is a compile-time filter}: a message above it is not merely
suppressed, it is not compiled into the image at all, format string included. That is why debug
messages can be left in the code for free, and why turning them on needs a rebuild.

\code{LOG_MODULE_REGISTER(monitor)}, with no level, uses \code{CONFIG_LOG_DEFAULT_LEVEL}, which is 3
unless something sets it. Zephyr’s own modules each have a Kconfig option for their level instead,
such as \code{CONFIG_I2C_LOG_LEVEL}, and an application can give its modules one the same way, with
a template from Zephyr’s logging Kconfig:

\begin{kconfigcode}
module = MONITOR
module-str = monitor
source "subsys/logging/Kconfig.template.log_config"
\end{kconfigcode}

\noindent That declares a choice of levels, set in \code{prj.conf} with a line such as
\code{CONFIG_MONITOR_LOG_LEVEL_DBG=y}, and the resulting number as \code{CONFIG_MONITOR_LOG_LEVEL}.
The module is registered with \code{LOG_MODULE_REGISTER(monitor, CONFIG_MONITOR_LOG_LEVEL)}.

\subsubsection{Deferred and immediate}
There are two ways a message can reach the console, and the choice matters on a microcontroller:

\begin{itemize}
  \item \textbf{Deferred} logging, the default on real hardware, only copies the message’s arguments into a
    buffer; a low-priority log thread formats and prints it later. Logging from an interrupt handler
    then costs microseconds, not the milliseconds a UART takes. The price is that a log message and a
    \code{printk()} right after it may appear in the opposite order, and that messages still in the
    buffer are lost if the system crashes.
  \item \textbf{Immediate} logging formats and prints in the caller, like \code{printk()}.
\end{itemize}

\textbf{\code{native_sim} logs in immediate mode by default} (\code{CONFIG_LOG_MODE_IMMEDIATE}),
because printing on a host costs nothing that matters and ordered output is easier to read. Keep in
mind that the same application on a board will be deferred unless told otherwise.

\subsection{Time, briefly}
Chapter~\ref{c3:ch} covers time properly. Two things are needed before then.

\textbf{Delays are written in milliseconds and kept in ticks.} \code{k_msleep(500)} and
\code{k_sleep(K_MSEC(500))} are the same call. The kernel converts the duration to \emph{ticks} of
its system clock, of which \code{native_sim} has \code{CONFIG_SYS_CLOCK_TICKS_PER_SEC=100}, one
every 10 ms. \code{k_uptime_get()} returns the milliseconds since boot.

\textbf{A sleep lasts at least as long as asked, and usually a little longer.} A thread that calls
\code{k_msleep(500)} at some moment within a tick cannot know how much of that tick is left, so the
kernel waits the 50 ticks the duration needs, plus the rest of the current one. That is why the
\file{lectures/L01/examples/hello} example wakes at 510 ms and 1020 ms, not at 500 and 1000. On
\code{native_sim} the numbers are the same on every run; on a board they would be the same too,
which is the point of a tick-based kernel.

\subsection{Summary}
\begin{itemize}
  \item Zephyr is a kernel, a driver model, subsystems, and a build that selects and configures all of them
    at compile time.
  \item \code{native_sim} builds Zephyr into a Linux program with simulated time: deterministic and fast,
    but a thread that never waits stops the world.
  \item A west workspace is a manifest and the repositories it pins. \code{west update} makes the workspace
    match the manifest.
  \item An application is \code{CMakeLists.txt}, \code{prj.conf} and sources. \code{find_package(Zephyr)} is
    where the build takes over.
  \item The build configures in order: board, devicetree, Kconfig, CMake. Then it compiles.
  \item A Kconfig value is the last one assigned, from the defaults through \code{prj.conf} to the command
    line, and an assignment that does not take aborts the build. \code{build/zephyr/.config} shows the
    result.
  \item Use \code{IS_ENABLED()} in C, and \code{*_ifdef()} in CMake. A misspelled \code{CONFIG_} macro in C
    is off, silently.
  \item Log through log modules. A module’s level decides at compile time which messages exist.
\end{itemize}

% Section 1.3, from lectures/L01/README.md: the objectives, the evaluation
% questions, and the next lecture.

\section{Review}
\label{c1:sec:review}

\needspace{6\baselineskip}
\subsection*{What you should be able to do now}
After this chapter, you should:

\begin{itemize}
  \item Be able to set up the workspace and build and run an application on \code{native_sim}.
  \item Understand what the west manifest pins, and where Zephyr lives in the workspace.
  \item Understand the three files every Zephyr application has, and what each one contributes.
  \item Understand the order in which a Zephyr build configures itself: board, devicetree, Kconfig, CMake.
  \item Be able to find the configuration a build ended up with, in \code{.config} and in menuconfig.
  \item Be able to add Kconfig options of your own and use them from C and CMake.
  \item Be able to log at the right level, and to change which levels a build keeps.
\end{itemize}

\needspace{6\baselineskip}
\subsection*{Questions to test yourself}
You should be able to explain:

\begin{enumerate}
  \item What does west do, and what does the manifest \code{west.yml} pin?
  \item What are the roles of \code{CMakeLists.txt}, \code{prj.conf} and devicetree in a Zephyr build?
  \item Where does a Kconfig option’s final value come from, and where can you read it after a build?
  \item Why does the build abort when \code{prj.conf} sets an option whose dependencies are not met?
  \item Why is a misspelled \code{CONFIG_} macro in C more dangerous than a misspelled option in
    \code{prj.conf}?
  \item What is the difference between \code{\#ifdef CONFIG_X} and \code{IS_ENABLED(CONFIG_X)}?
  \item Why does \code{k_msleep(1000)} on \code{native_sim} wake up after 1010 ms?
  \item What does a log module’s level control, and why is a filtered-out message not in the binary?
\end{enumerate}

\needspace{6\baselineskip}
\subsection*{Looking ahead}
Chapter~\ref{c2:ch} is about devicetree:

\begin{itemize}
  \item How a board and its peripherals are described: nodes, properties, labels and aliases.
  \item Bindings, which say what a node’s properties mean.
  \item Overlays, which change a board’s devicetree for one application.
  \item Getting devices and properties from C at compile time.
\end{itemize}

You will blink an LED and read a button through Zephyr’s emulated GPIO controller.

% Section 1.4, from lectures/L01/appendix/c_exercises.md.

\section{Exercises}
\label{c1:sec:exercises}\label{c1:app:c}
These exercises reinforce the concepts from \secref{c1:app:b}.

\begin{aside}
\textbf{How to check your work.} Write each exercise set as its own application under
\file{lectures/L01/exercises}, with the file names given there, and run \code{make test} from the
repository root. A set’s tests switch on as soon as its files exist, and a set you have not started
is reported as \code{SKIP}. Each application grows through its set, and the tests check where it
ends up. Programs run on \code{native_sim} in simulated time, so their output, timestamps included,
is compared with the expected output line by line.

The solutions are published in the companion repository, under
\file{lectures/L01/appendix/solutions}, a few days after the lecture. Try each exercise yourself
before you read them.
\end{aside}

Remember \code{source env.sh} in every new terminal. Build and run an exercise from the repository
root:

\begin{shell}
west build -p -b native_sim/native/64 lectures/L01/exercises/exercise1
west build -t run
\end{shell}

% ---------------------------------------------------------------------
\exerciseset{A first application}
\label{c1:set:1}
Create a directory named \code{exercise1} to work in for this exercise set.

\exercise{The three files}{Code}
\label{c1:ex:1.1}
In this exercise you will create the smallest Zephyr application that builds.

\task{Tasks}
Create these files in \code{exercise1}:

\begin{itemize}
  \item \code{CMakeLists.txt}: require CMake 3.20, find the Zephyr package, name the project
    \code{exercise1}, and add \code{src/main.c} to the \code{app} target.
  \item \code{prj.conf}: an empty file, or one holding only a comment.
  \item \code{src/main.c}: a \code{main()} that prints \code{Hello, Zephyr!} with \code{printk()} and
    returns \code{0}.
\end{itemize}

Build and run the application.

\task{Reflection}
\begin{itemize}
  \item What happens if you delete \code{prj.conf} and build again with \code{-p}?
  \item The program prints its greeting and then keeps running. Why does it not exit when \code{main()}
    returns? How do you stop it?
  \item Where in \code{build/} is the program you ran?
\end{itemize}

\exercise{The uptime}{Code}
\label{c1:ex:1.2}

\task{Tasks}
Extend \code{main()} so that after the greeting it prints the uptime three times, with
\code{k_msleep(1000)} after each print:

\begin{console}
Hello, Zephyr!
Uptime: 0 ms
Uptime: 1010 ms
Uptime: 2020 ms
\end{console}

Use \code{k_uptime_get()}, which returns an \code{int64_t}, printed with \code{\%lld}.

Run the program with \code{./build/zephyr/zephyr.exe -stop_at=5}, and then again with
\code{-stop_at=5 -no-rt}.

\task{Reflection}
\begin{itemize}
  \item Why is the second uptime 1010 ms and not 1000 ms? Find the option in \code{build/zephyr/.config}
    that decides it.
  \item What is the difference between the two runs? Is the output different?
  \item Would the uptimes be the same if the host were busy compiling something else at the same time? Why?
\end{itemize}

% ---------------------------------------------------------------------
\exerciseset{Logging}
\label{c1:set:2}
Create an application named \code{exercise2} to work in for this exercise set, with the same three
files as \code{exercise1}.

\exercise{A temperature monitor}{Code}
\label{c1:ex:2.1}
In this exercise you will replace \code{printk()} with the logging subsystem.

\task{Tasks}
\begin{itemize}
  \item Enable logging in \code{prj.conf}.
  \item Register a log module named \code{monitor} \textbf{without} giving a level:
    \code{LOG_MODULE_REGISTER(monitor);}.
  \item In \code{main.c}, define the readings as an array:
\begin{ccode}
static const int readings[] = {21, 24, 28, 33, 36, 29};
\end{ccode}
  \item In \code{main()}, go through the readings, and for each one log:
\begin{itemize}
  \item \code{Temperature <t> C is critical} as an error, if it is 35 or more.
  \item \code{Temperature <t> C is high} as a warning, if it is 30 or more.
  \item \code{Temperature <t> C} as information, otherwise.
\end{itemize}
  \item Sleep 500 ms after each reading.
\end{itemize}

The output must be:

\begin{console}
[00:00:00.000,000] <inf> monitor: Temperature 21 C
[00:00:00.510,000] <inf> monitor: Temperature 24 C
[00:00:01.020,000] <inf> monitor: Temperature 28 C
[00:00:01.530,000] <wrn> monitor: Temperature 33 C is high
[00:00:02.040,000] <err> monitor: Temperature 36 C is critical
[00:00:02.550,000] <inf> monitor: Temperature 29 C
\end{console}

\note[Hint] \code{ARRAY_SIZE(readings)}, from \code{<zephyr/sys/util.h>}, which
\code{<zephyr/kernel.h>} includes, gives the number of elements.

\exercise{Only what matters}{Code}
\label{c1:ex:2.2}

\task{Tasks}
Without changing any file, build the application so that it logs only warnings and errors, by
setting the default log level on the command line:

\begin{shell}
west build -p -b native_sim/native/64 lectures/L01/exercises/exercise2 -- -DCONFIG_LOG_DEFAULT_LEVEL=2
\end{shell}

\noindent The output must now be only:

\begin{console}
[00:00:01.530,000] <wrn> monitor: Temperature 33 C is high
[00:00:02.040,000] <err> monitor: Temperature 36 C is critical
\end{console}

\task{Reflection}
\begin{itemize}
  \item Why did this change the monitor’s output, when the command line said nothing about \code{monitor}?
  \item Would it still work if the module were registered as \code{LOG_MODULE_REGISTER(monitor,
    LOG_LEVEL_INF)}?
  \item Are the information messages still in the program, only not printed? Explain how you know.
  \item After this build, run \code{west build -t run} a few times, and then build again \textbf{without}
    \code{-p} and without the \code{-D} option. Which log level does that build use? Look for the answer
    in \code{build/zephyr/.config}, and explain it.
\end{itemize}

% ---------------------------------------------------------------------
\exerciseset{Options of your own}
\label{c1:set:3}
Create an application named \code{exercise3} to work in for this exercise set.

\exercise{The threshold}{Code}
\label{c1:ex:3.1}
In this exercise you will give the application Kconfig options of its own.

\task{Tasks}
\begin{itemize}
  \item Create a \code{Kconfig} file in \code{exercise3} with a menu named \code{Application}, holding:
\begin{itemize}
  \item \code{APP_HIGH_C}, an \code{int} with the prompt \code{Temperature at which a reading is high, in
    degrees Celsius}, a range from -40 to 125, and a default of 30.
\end{itemize}
  \item End the file so that Zephyr’s own options are still available.
  \item In \code{prj.conf}, set \code{APP_HIGH_C} to 25.
  \item In \code{src/main.c}, go through these readings, printing each with \code{printk()}, and mark those
    at or above \code{CONFIG_APP_HIGH_C} with \code{ (high)}:
\begin{ccode}
static const int readings[] = {-18, 21, 24, 28, 33};
\end{ccode}
\end{itemize}

The output must be:

\begin{console}
-18 C
21 C
24 C
28 C (high)
33 C (high)
\end{console}

\task{Reflection}
\begin{itemize}
  \item Find \code{APP_HIGH_C} in \code{west build -t menuconfig}. What does \code{?} show about it?
  \item Set \code{CONFIG_APP_HIGH_C=200} in \code{prj.conf} and build. What happens, and why is that better
    than the alternative?
  \item Remove the last line of your \code{Kconfig} file and build with \code{-p}. What happens?
\end{itemize}

\exercise{Fahrenheit, when asked for}{Code}
\label{c1:ex:3.2}

\task{Tasks}
\begin{itemize}
  \item Add a \code{bool} option \code{APP_FAHRENHEIT} to the \code{Application} menu, with the prompt
    \code{Print temperatures in degrees Fahrenheit}.
  \item Create \code{src/convert.h} and \code{src/convert.c} with a function:
\begin{ccode}
int celsius_to_fahrenheit(int celsius);
\end{ccode}
which returns the temperature in degrees Fahrenheit, \textbf{rounded to the nearest degree}, using
only integer arithmetic. For example, 37 °C is 98.6 °F, which rounds to 99, and -18 °C is -0.4 °F,
which rounds to 0.
  \item In \code{CMakeLists.txt}, compile \code{src/convert.c} \textbf{only} when
    \code{CONFIG_APP_FAHRENHEIT} is enabled.
  \item In \code{main()}, when the option is enabled, print each reading in Fahrenheit, as \code{<t> F},
    still marked \code{ (high)} by its value in Celsius. Use \code{IS_ENABLED()}, not \code{\#ifdef}.
  \item Create a file \code{fahrenheit.conf} in \code{exercise3} that enables the option.
\end{itemize}

Built with \code{-DEXTRA_CONF_FILE=fahrenheit.conf}, the output must be:

\begin{console}
0 F
70 F
75 F
82 F (high)
91 F (high)
\end{console}

\noindent Built without it, the output must be unchanged from Exercise~\ref{c1:ex:3.1}.

\note[Hint] Work in tenths of a degree: 37 °C is 986 tenths of a degree Fahrenheit. C’s integer
division truncates towards zero, so \code{(tenths + 5) / 10} rounds a positive number of tenths to
the nearest degree, but not a negative one.

\task{Reflection}
\begin{itemize}
  \item Without \code{fahrenheit.conf}, \code{convert.c} is not compiled, yet \code{main.c} still calls
    \code{celsius_to_fahrenheit()}. Why does the program link?
  \item Would it still link with \code{\#ifdef CONFIG_APP_FAHRENHEIT} around the call instead? Which of the
    two would you rather maintain, and why?
  \item Change \code{IS_ENABLED(CONFIG_APP_FAHRENHEIT)} to \code{IS_ENABLED(CONFIG_APP_FARENHEIT)} and build
    with \code{fahrenheit.conf}. What happens? What would have caught it?
\end{itemize}

% ---------------------------------------------------------------------
\exerciseset{Reading the workspace and the build}
\label{c1:set:4}
No code in this set: answer the questions by looking. Use the build of any exercise above.

\exercise[fillaccent2]{The workspace}{Reflection}
\label{c1:ex:4.1}

\task{Tasks}
Answer, using west’s commands rather than by opening files:

\begin{itemize}
  \item Which revision of Zephyr is the workspace using, and where is it checked out?
  \item Which configuration value makes \code{west build} work without \code{-b}?
\end{itemize}

\exercise[fillaccent2]{The build}{Reflection}
\label{c1:ex:4.2}

\task{Tasks}
Answer, by looking in the build directory:

\begin{itemize}
  \item Which log mode (\code{CONFIG_LOG_MODE_...}) did the \code{exercise2} build use? Why that one?
  \item What is the value of \code{CONFIG_BOARD}, and what is the value of \code{CONFIG_BOARD_TARGET}? Why
    are they different?
  \item How many \code{\#define} lines does \code{autoconf.h} hold? Does it contain the options that are
    \code{n}?
  \item Find the definition of \code{k_msleep()} with \code{grep -n -A3 "k_msleep(int32_t ms)"
    deps/zephyr/include/zephyr/kernel.h}. What does it call, and what does \code{static inline} mean for
    whether \code{k_msleep} exists as a function in your program?
\end{itemize}

